% Nouveau EGA V (cote 25): archive cote 25, site page 110 (PDF page 111).
% typescript. Draft compared: none.
% AI-transcribed by Claude Opus 5.5 (claude-opus-5-5, adaptive thinking, effort high), Message Batches API, batch ega25b (msgbatch_01A7n5YuXubUPNK1zsepi8RG), 2026-11-13; not yet checked by eye.
\scanpage{25/110}
\mspage{13}

Nous nous proposons de déterminer le sous-groupe Div$^{\mathrm{L}}$(X/S) de Div$^{+}$(X/S)
formé des diviseurs relatifs positifs tels que $\underline{O}_X$(D) soit isomorphe à
un Module de la forme $\underline{L}\otimes_{\underline{O}_X}\underline{M}$, où $\underline{M}$ est un Module inv sur S (dépendant de
D). Nous supposerons qu'on a
\[ f_{*}(\underline{O}_X)\overset{\sim}{\longleftarrow}\underline{O}_S \quad , \]
ce qui implique que ci-dessus $\underline{M}$ est déterminé à isom unique près en termes
de D, comme
\[ \underline{M}\simeq f_{*}(\underline{L}^{-1}\otimes\underline{O}_X(D)) \quad . \]
La donnée d'un D correspondant à un $\underline{M}$ donné revient donc à la donnée
\inserted[typed]{\xout{$\underline{\phi}$x}}
d'une section\xout{s} $\phi$ de $\underline{L}\otimes\underline{M}$ transversalement régulière. Or on a vu au
Chap III par 7 \xout{f$_*$($\underline{L}\otimes\underline{M}$)xxxxf$_*$($\underline{L}$)$\otimes$Mxx} qu'il \xout{xixx} existe un Module Q de
présentation finie sur S, dont la formation commute d'ailleurs à tout
changement de base, et un isomorphisme de foncteurs en \unsure{$\underline{N}$|$\underline{G}$} (Module qu coh
variable sur S)
\[ f_{*}(\underline{L}\otimes\underline{G})\simeq \mathrm{Hom}(\underline{Q},G) \quad . \]
(A vrai dire, dans loc cit on a supposé S loc noethérien, mais on se
débarasse de cette hypothèse de façon évidente, par le pro\xout{x}cédé brévété,
compte tenu de la commutation de Q au changement de\xout{x} base). En particulier
la donnée de $\phi$ équivaut à \xout{xx} la donnée d'un homomorphisme
\[ \psi : Q \longrightarrow M \ . \]
Une condition nécessaire pour que $\phi$ soit transversalement régulière, et
qui est suffisante lorsque les fibres de X sont intégres, est que $\phi$ soit
$\neq 0$ \xout{xix} fibre par fibre, ce qui en termes de $\psi$ signifie simplement que
$\psi$ est surjectif, donc que \xout{Mxxxx} $\psi$ correspond à une section du fibré pro-%
jectif $\underline{P}(\underline{Q})$ sur S. On obtient ainsi la

\uline{Proposition} 2.1\unsure{1}. \inserted[typed]{Avec les notations} \xout{Sous les conditions} précédentes,\xout{xxx X/S} Div$^{\mathrm{L}}$(X/S) est
en correspondance biunivoque canonique avec l'ensemble des \xout{point} sections
de $\underline{P}(\underline{Q})$ sur S, correspondant à des Modules quotients $\underline{M}$ de $\underline{Q}$ tel\unsure{\xout{s}|s} que

\mspage{110}
